\scanpage{70}
% loc: 1–2, diagonally struck
\deleted{utilisant la formule}
\[
 \deleted{(a\overset{e}{\otimes}\underbrace{b)\circ(c}\overset{e}{\otimes}d)
 \simeq c\overset{e}{\otimes}b\,(d\wedge a/e),}
\]

% loc: 3–5, same cancellation
\deleted{on trouve}
\[
\begin{aligned}
 &\deleted{uvw=(y\overset{e}{\otimes}x\,\underbrace{(-x\wedge y/e)}_{-1})\,z\overset{e}{\otimes}z}\\
 &\deleted{\phantom{uvw}=-z\otimes x\,\underbrace{(z\wedge y/e)}_{-1}=z\overset{e}{\otimes}x}
\end{aligned}
\]

% loc: 6–8
Il existe donc un unique $\rho\in\mathrm{Gl}(E)$ tel que
\[
 \rho x=y,\ \rho y=z,\ \rho z=x
\]

% loc: 9–13
et ce $\rho$ est d'ordre 3. Il ne change pas en remplaçant $\unsure{\{x,y,z\}|(x,y,z)}$ par $(-x,-y,-z)$, il ne dépend que du triplet $\{\dot{x},\dot{y},\dot{z}\}$ et d'une orientation de ce triplet. Changer l'orientation, revient à changer $\rho$ en $\rho^{-1}$

% loc: 14–15
\unsure{Or on aura|On en aura|On a aussi}
\[
 \det\rho=1,\ \operatorname{Tr}\rho=1,\ \inserted{\text{donc}}\quad
 \rho^{-1}=\deleted{\unsure{\rho}}1-\rho=\rho'\quad(\text{conj. de }\rho)
\]

% loc: 16–17
Considérons
\[
 \unsure{\sigma_f|\sigma|\tau} =x\overset{e}{\otimes}x+y\overset{e}{\otimes}y+z\overset{e}{\otimes}z
 \qquad\text{On a }\left|\begin{aligned}
 \operatorname{Tr}\sigma&=0\\
 \det\sigma&=\deleted{4}3.
 \end{aligned}\right.\quad
 \text{(par formule d'additivité du dét\ldots)}
\]

% loc: 18–20
c'est un endom qui commute à $\rho$, donc de la forme
\[
 \sigma=a\rho+b
\]

% loc: 21–24, bottom
\deleted{\unsure{où}} d'où, en prenant les traces
\[
\begin{aligned}
 &\operatorname{Tr}\sigma.=a+2b\qquad\text{donc }\underset{\deleted{\unsure{b=-a}}}{a=-2b},\ \text{donc}\\
 &\sigma=\deleted{\unsure{a(-2b|a\rho+2b}}\,b(-2\rho+1)\\
 &\det\sigma=b^2\,\underline{\det(1-2\rho)},\qquad \text{or }\det(1-2\rho)=3
\end{aligned}
\]
% archive sheet number: 69

